\section{Conclusão}

Apesar de atingir o objetivo original, o presente estudo deixa a desejar quanto
a seus resultados no que diz respeito ao catálogo final gerado. Como pode-se
observar na seção de resultados, os métodos empregados foram capazes de detectar
apenas padrões relativamente simples e óbvios.

Apesar de representativos de diferentes formas de minerar padrões de uso de APIs
na literatura, o estudo foi extremamente restrito devido ao escopo e escolha de
método de análise sintática dos arquivos. Como mencionado anteriormente, a
escolha considerar apenas métodos diretamente chamados na mesma função pode
levar a perda de padrões de uso reais. Ademais, outros metodologias que se
utilizavam de entradas mais complexas ( como rastros de execução do código por
exemplo ) precisaram ser desconsideradas do presente estudo por não haverem os
dados necessários para aplicá-las.

Dessa forma, um futuro estudo no mesmo tópico poderia ter mais sucesso buscando
formas mais rebuscadas de avaliar quais funções são chamadas em conjunto,
permitindo uma visão mais completa do que é realmente feito para os métodos já
testados. Tal também permitiria a aplicação de outras formas de minerar padrões
de uso mais complexas e diversas.
